\section{Visitation and Definitive Constitutions, 2014--2016}

Experimental constitutional approval in 2008 created a scheduled need for review. In 2014 an ICKSP news page announced that the Institute was undergoing an ``ordinary apostolic visitation'' ordered by the Pontifical Commission \emph{Ecclesia Dei}. It named Archbishop François Bacqué, a retired apostolic nuncio, as visitor and Dominican Luc-Thomas Somme as co-visitor. The page framed visitation as normal maternal oversight five years after pontifical-right erection.

That public notice establishes a mandate and the Institute's characterization. It does not publish the Commission's original mandate, itinerary, testimony, findings, recommendations, or final report. ``Ordinary'' means the process was presented as a normal supervisory act, not that a favorable result is known in advance. A visitation is not inherently penal, but neither is its existence a certificate that every practice was examined and approved.

\statusframe{ICKSP, ``Visite canonique,'' institutional news page in 2014.}{PCED ordered an ordinary apostolic visitation and appointed Bacqué with Somme to carry it out.}{The original mandate and report were not located. The Institute's benign framing is attributed and no dispositive outcome is inferred.}

\subsection{The 2016 institutional notice}

An ICKSP page records that the Holy See definitively approved the Institute's constitutions on 29 January 2016, the feast of Saint Francis de Sales. It reports a thanksgiving \emph{Te Deum} and says Bacqué, identified as the former apostolic visitor, would celebrate a pontifical Mass at Gricigliano on 2 February. In his signed sermon for that Mass, also hosted by the Institute, Bacqué said that the Holy See---specifically the Congregation for the Doctrine of the Faith incorporating PCED---had definitively approved the constitutions. This is a direct contemporary witness by the former visitor to the approval and competent Roman office. It does not print the act, its provisions, or the 29 January date supplied by the institutional notice.

The paired notice and sermon strongly support movement from the five-year experimental code to definitive proper law. They do not show which provisions changed, whether every recommendation was adopted, or whether definitive approval should be described as the visitation's ``verdict.'' Bacqué confirms the later Roman act but does not disclose the visitation's findings or causal relation to approval.

Definitive approval also does not mean immutability. Proper law can be amended through competent procedures, subordinate directories can develop, and universal law remains controlling. ``Definitive'' distinguishes the constitutional approval from the experimental term; it does not promise that no later visitation, reform, or interpretation will occur.

\subsection{Government after the settlement}

The 2008 decree named Wach for a renewable six-year term. A 2020 institutional announcement says a General Chapter met at Gricigliano under article 20 of the constitutions, in the presence of Monsignor Patrick Incorvaja, whom the notice associates with the Congregation for the Doctrine of the Faith, and elected Wach superior general for another six years. That public election record illustrates ordinary constitutional government after definitive approval, though the complete article, voting procedure, and chapter acts remain unavailable.

The concentration of founding and later government in Wach is historically significant. It supplies continuity of charism and networks over decades. It also makes independent supervision and constitutional mechanisms important: institutional memory, formation, appointments, and external representation can become closely identified with one founder. The historical fact is continuity; whether particular exercises of authority were prudent belongs to evidence and competent review, not inference from tenure alone.

\claimcheck{``The 2014 visitation cleared the Institute and led to approval in 2016.''}{The public record shows an ordinary visitation and Bacqué's direct contemporary confirmation of later definitive constitutional approval. Without the mandate, report, approval act, or a causal statement, the stronger claim still invents both a verdict and a documented causal chain.}
